\chapter{Contexto, alcance y supuestos}

\section{Contexto empresarial}

RedNova GT desarrolla una plataforma SaaS para inventario, facturación y análisis de ventas. El equipo de Ventas consulta demostraciones y oportunidades comerciales; TI administra sistemas, identidades y soporte; el Data Center aloja servicios internos y bases de datos; y los visitantes requieren conectividad aislada. Aunque la organización es ficticia, la separación responde a un patrón empresarial real: distintos grupos tienen necesidades y niveles de confianza distintos, por lo que compartir una red plana aumentaría tanto el dominio de falla como la superficie de exposición.

La fase no intenta construir una plataforma productiva de alta disponibilidad. Su propósito es producir un prototipo verificable que demuestre diseño, subneteo, configuración base y conectividad interna. La solución se concentra en una sola región y una sola zona de disponibilidad para reducir complejidad y costo. Una versión productiva debería duplicar los segmentos críticos en al menos dos zonas, incorporar monitoreo, registro centralizado, respaldo, balanceo y mecanismos administrados de acceso.

\section{Objetivos}

\subsection{Objetivo general}

Diseñar y prototipar en AWS una red empresarial segmentada para RedNova GT, demostrando mediante cálculos y pruebas que la asignación de direcciones, las rutas y los controles de acceso cumplen las necesidades declaradas.

\subsection{Objetivos específicos}

\begin{itemize}
  \item Comparar proveedores de nube con criterios técnicos, académicos y económicos.
  \item Traducir cantidades de usuarios y servidores a prefijos CIDR válidos para AWS.
  \item reservar espacio suficiente para crecimiento sin exceder el bloque principal.
  \item Separar segmentos públicos y privados mediante tablas de rutas explícitas.
  \item Aplicar mínimo privilegio con grupos de seguridad orientados a función.
  \item Automatizar la infraestructura para reducir errores manuales y permitir repetición.
  \item Demostrar conectividad ICMP entre dos instancias privadas dentro de Ventas.
  \item Documentar costos, riesgos, evidencias y el procedimiento de destrucción.
\end{itemize}

\section{Alcance incluido}

\begin{tabularx}{\textwidth}{p{4.4cm}Y}
\toprule
\textbf{Componente} & \textbf{Cobertura de esta fase} \\
\midrule
Direccionamiento & Una VPC IPv4 privada, cinco subredes implementables y un bloque de reserva. \\
Enrutamiento & Tabla pública con ruta por Internet Gateway y tabla privada sin salida directa. \\
Seguridad & Dos grupos de seguridad, restricción SSH por \code{/32} e ICMP interno controlado. \\
Cómputo & Un bastión temporal y dos instancias pequeñas de prueba cuando se habilita \code{enable\_compute}. \\
Automatización & Código Terraform, variables, salidas, etiquetas y scripts de validación. \\
Pruebas & Capacidad CIDR, no solapamiento, validación Terraform y ping mediante direcciones privadas. \\
Documentación & Informe, diagramas, evidencias y README de la fase. \\
\bottomrule
\end{tabularx}

\section{Fuera de alcance}

Quedan fuera de esta fase: arquitectura multirregión, alta disponibilidad en múltiples zonas, VPN hacia oficinas físicas, directorio corporativo, balanceadores, bases de datos administradas, inspección profunda de paquetes, IPv6, observabilidad completa, pruebas de rendimiento y recuperación ante desastres. Estos elementos aparecen como recomendaciones de evolución, no como componentes supuestamente implementados.

\section{Supuestos y restricciones}

\begin{longtable}{p{4.0cm}p{8.2cm}p{2.1cm}}
\toprule
\textbf{Supuesto} & \textbf{Justificación e impacto} & \textbf{Estado} \\
\midrule
\endfirsthead
\toprule
\textbf{Supuesto} & \textbf{Justificación e impacto} & \textbf{Estado} \\
\midrule
\endhead
30 visitantes concurrentes & El enunciado no cuantifica Visitas. Se adopta una carga razonable que exige \code{/26} en AWS y deja margen. & Explícito \\
25 dispositivos en Ventas & Se interpreta una dirección primaria por colaborador, más crecimiento operativo separado. & Dado \\
15 dispositivos en TI & Se interpreta una dirección primaria por colaborador; \code{/27} deja doce libres. & Dado \\
5 servidores lógicos & La cifra dimensiona Data Center; no obliga a ejecutar cinco EC2 durante el prototipo. & Dado \\
Una zona de disponibilidad & Es suficiente para demostrar segmentación, pero no ofrece tolerancia a la falla de zona. & Limitación \\
Acceso administrativo conocido & El equipo debe conocer su IPv4 pública y expresarla como \code{/32}. & Operativo \\
Cuenta con permisos mínimos & Se requieren permisos para VPC, EC2, grupos de seguridad, llaves y etiquetas. & Externo \\
\bottomrule
\end{longtable}

\begin{warningbox}[Restricción económica]
El diseño evita NAT Gateway y direcciones públicas innecesarias. Sin embargo, AWS cobra por direcciones IPv4 públicas asociadas y ociosas \cite{aws-vpc-pricing}. El bastión sólo debe existir durante la prueba y su dirección debe liberarse al destruir el prototipo.
\end{warningbox}
